\label{sec:numerical-examples}

This section turns from problem definition to computational behavior. The
benchmark families in \Cref{sec:benchmarks} established the governing
equations, discrete energies, and representative states, and
\Cref{sec:validation} separated the validation and external-comparison
evidence. Here we focus on parallel timings, scaling trends, solver behavior,
and the comparative cost of alternative derivative or discretization choices.
Evidence from older timing studies is retained when it is still informative for
workflow viability, but it is labeled as timing evidence rather than merged into
the validation story.

\subsection{\texorpdfstring{$p$-Laplace}{p-Laplace}}
\label{sec:num-plaplace}

The $p$-Laplace problem definition and discrete energy are given in
\Cref{sec:bench-plaplace,eq:plaplace-discrete}. In this clean scalar setting,
the numerical story is easiest to interpret: exact element Hessians remain
effective in the reported distributed runs, while local-SFD reaches nearly
identical discrete energies at higher wall time.

\begin{figure}[H]
  \centering
  \includegraphics[width=\PaperFigureMedium]{plaplace_scaling.pdf}
  \caption{$p$-Laplace strong scaling on the finest maintained grid. All curves
  correspond to the same discrete problem from \cref{eq:plaplace-discrete}.}
  \label{fig:plaplace-results}
\end{figure}

\subsection{Ginzburg--Landau}
\label{sec:num-gl}

The Ginzburg--Landau family from \Cref{sec:bench-gl,eq:gl-discrete} shows the
same high-level derivative story under a more indefinite linearization. The
important computational message is not a dramatic separation between exact and
approximate second-order routes, but that the maintained distributed solver
remains convergent on this nonconvex scalar energy with the same sparse-infrastructure
contract.

\begin{figure}[H]
  \centering
  \includegraphics[width=\PaperFigureMedium]{ginzburg_landau_scaling.pdf}
  \caption{Ginzburg--Landau strong scaling for the maintained fine-grid case
  governed by \cref{eq:gl-discrete}.}
  \label{fig:gl-results}
\end{figure}

\subsection{Hyperelasticity}
\label{sec:num-hyperelasticity}

The hyperelasticity problem definition, boundary data, and terminal energies
are given in \Cref{sec:bench-he,eq:hyperelasticity-discrete}. Computationally,
this is the maintained benchmark in which globalization matters more than
derivative-route choice: the distributed implementation completes the full load
path and preserves the terminal energy trend from \Cref{sec:bench-he}. The
external JAX-FEM companion comparison is reported separately in
\Cref{sec:val-hyperelasticity}; the figure below reports only distributed
scaling for the repository load path.

\begin{figure}[H]
  \centering
  \includegraphics[width=\PaperFigureMedium]{hyperelasticity_scaling.pdf}
  \caption{Hyperelasticity strong scaling for the 24-step load path defined by
  \cref{eq:hyperelasticity-discrete}.}
  \label{fig:hyperelasticity-results}
\end{figure}

The larger fixed-work Karolina experiment isolates the first level-5 load
step after the rank-local HyperElasticity implementation. The Karolina CPU
nodes used here are IT4I standard CPU nodes with two AMD EPYC 7H12 64-core
processors and 256 GB RAM \citep{it4i-karolina-compute-nodes}; on the tested
nodes Slurm exposes the 128 cores as eight NUMA domains with 16 cores each.
The runs use one MPI rank per core, contiguous CPU-ID binding, and local memory
binding. The mesh is generated procedurally from the uniform beam hierarchy,
and the element path assembles from rank-local elements with local PETSc COO
preallocation and point-to-point overlap exchange.

The nonlinear solve is the same trust-region Newton path used for the
benchmark in \cref{sec:bench-he}. The linear subproblems use STCG with a
same-hierarchy PMG preconditioner, Chebyshev/Jacobi smoothing, and coarsest
level $L_2$. The coarse solve is PETSc redundant LU with MUMPS: the partial
one-node rows use one redundant group over the active ranks, while the
multi-node rows use one 128-rank redundant group per node.
\Cref{fig:hyperelasticity-karolina-pmg-scaling,tab:hyperelasticity-karolina-pmg-scaling}
report the application-reported solver total time, not Slurm elapsed time.
All rows reach the same first-step energy to the precision shown in the raw
outputs, with 21 Newton iterations and 43 Krylov iterations.

\begin{figure}[H]
  \centering
  \includegraphics[width=\PaperFigureMedium]{hyperelasticity_karolina_pmg_scaling.pdf}
  \caption{Karolina level-5 HyperElasticity first-step fixed-work scaling with
  the PMG coarse level fixed at $L_2$. Multi-node annotations show node counts.}
  \label{fig:hyperelasticity-karolina-pmg-scaling}
\end{figure}

\begin{table}[H]
  \centering
  \begin{tabular*}{\textwidth}{@{}c@{\extracolsep{\fill}} c c c c c c c@{}}
\toprule
Ranks & Nodes & Coarse groups & Ranks/group & Solver total [s] & First step [s] & Newton iters & Krylov iters \\
\midrule
16 & 1 & 1 & 16 & \num{419} & \num{367} & 21 & 43 \\
32 & 1 & 1 & 32 & \num{317} & \num{280} & 21 & 43 \\
64 & 1 & 1 & 64 & \num{148} & \num{132} & 21 & 43 \\
128 & 1 & 1 & 128 & \num{67.4} & \num{58.4} & 21 & 43 \\
256 & 2 & 2 & 128 & \num{36.4} & \num{30.7} & 21 & 43 \\
512 & 4 & 4 & 128 & \num{21.7} & \num{17.7} & 21 & 43 \\
1024 & 8 & 8 & 128 & \num{15.7} & \num{12.3} & 21 & 43 \\
\bottomrule
\end{tabular*}

  \caption{Karolina level-5 HyperElasticity first-step timings used in
  \cref{fig:hyperelasticity-karolina-pmg-scaling}.}
  \label{tab:hyperelasticity-karolina-pmg-scaling}
\end{table}

\subsection{Plasticity2D}
\label{sec:num-plasticity2d}

The Plasticity2D model and Davis-B reduction are defined in
\Cref{sec:bench-plasticity2d,eq:plasticity2d-discrete,eq:plasticity2d-davis}.
For this family, the most informative computational evidence is not a single
strong-scaling curve but the solver-engineering material collected in
\Cref{sec:solver-evidence}: same-mesh PMG choices, fixed-work runs, and source
continuation behavior. We therefore keep the within-benchmark resolution trend
in \Cref{fig:plasticity2d-energy-levels} and defer the supporting solver data
to the later evidence section.

\subsection{Plasticity3D}
\label{sec:num-plasticity3d}

The governing problem, Davis-B reduction, scalar constitutive potential, and
deviatoric-strain norm are defined in
\Cref{sec:bench-plasticity3d,eq:plasticity3d-discrete,eq:plasticity3d-davis,eq:plasticity3d-potential,eq:plasticity3d-devnorm}.
The validation limits for this endpoint surrogate are reported in
\Cref{sec:val-plasticity3d}. With that interpretation fixed, this section asks
two performance questions: which second-order route is preferred on the locked
flagship case, and how the maintained surrogate behaves across degree,
resolution, and rank count.

\paragraph{Derivative-route ablation.}
The first question is which second-order route has the most useful
cost-quality profile for the promoted maintained workflow. On the locked
$P_4(L_1)$ case at $\lambda_{\mathrm{sr}}=1.5$ and 8 MPI ranks, all three
routes converge to essentially the same terminal state, with identical Newton
and total Krylov counts. The differences are therefore cost differences, not
state-quality differences: constitutive AD has the lowest median wall time on
this locked case, with median wall time \SI{222}{s} and median solve time
\SI{191}{s}, compared with \SI{288}{s}/\SI{255}{s} for element AD
and \SI{372}{s}/\SI{253}{s} for colored SFD.

\begin{figure}[H]
  \centering
  \includegraphics[width=\PaperFigureMedium]{plasticity3d_derivative_ablation_bars.pdf}
  \caption{Locked Plasticity3D derivative-route ablation on
  $P_4(L_1)$, $\lambda_{\mathrm{sr}}=1.5$, and 8 MPI ranks.}
  \label{fig:plasticity3d-derivative-ablation}
\end{figure}

\begin{table}[H]
  \centering
  \begin{tabularx}{\textwidth}{@{}>{\hsize=0.88\hsize\linewidth=\hsize\RaggedRight\arraybackslash}Xr r r r r r r@{}}
\toprule
Route & Wall [s] & Solve [s] & Newton & Linear & Energy & $\omega$ & $u_{\max}$ \\
\midrule
Element AD & 288.021 & 254.879 & 8.00 & 376.0 & -3010860.766733 & 6027722.264408 & 0.731848 \\
Constitutive AD & 221.689 & 190.503 & 8.00 & 376.0 & -3010860.766439 & 6027722.229618 & 0.731848 \\
Colored SFD & 371.993 & 252.814 & 8.00 & 376.0 & -3010860.766733 & 6027722.264408 & 0.731848 \\
\bottomrule
\end{tabularx}

  \caption{Plasticity3D derivative-route ablation on the locked
  $P_4(L_1)$, $\lambda_{\mathrm{sr}}=1.5$, 8-rank case.}
  \label{tab:plasticity3d-derivative-ablation}
\end{table}

\paragraph{Maintained discretization and historical scaling.}
The remaining figures in this subsection answer a different question: how the
maintained surrogate behaves across degree/resolution choices, and how the
recommended maintained stack scales on the separate promoted
$P_4(L_2)$ timing study.

\begin{figure}[H]
  \centering
  \begin{subfigure}[t]{0.46\PaperFigureFull}
    \centering
    \includegraphics[width=\linewidth]{plasticity3d_degree_energy_all_dofs.pdf}
    \caption{All degrees: final energy versus free DOFs.}
  \end{subfigure}
  \hfill
  \begin{subfigure}[t]{0.46\PaperFigureFull}
    \centering
    \includegraphics[width=\linewidth]{plasticity3d_degree_energy_all_time.pdf}
    \caption{All degrees: final energy versus total wall time.}
  \end{subfigure}

  \vspace{0.45em}

  \begin{subfigure}[t]{0.46\PaperFigureFull}
    \centering
    \includegraphics[width=\linewidth]{plasticity3d_degree_energy_zoom_dofs.pdf}
    \caption{$P2/P4$ zoom: final energy versus free DOFs.}
  \end{subfigure}
  \hfill
  \begin{subfigure}[t]{0.46\PaperFigureFull}
    \centering
    \includegraphics[width=\linewidth]{plasticity3d_degree_energy_zoom_time.pdf}
    \caption{$P2/P4$ zoom: final energy versus total wall time.}
  \end{subfigure}
  \caption{Corrected glued-bottom maintained Plasticity3D
  degree-versus-resolution study at $\lambda_{\mathrm{sr}} = 1.55$. The top
  row shows the full $P_1/P_2/P_4$ comparison, while the bottom row zooms to
  $P_2/P_4$ so the near-overlap hidden by the coarse $P_1$ points becomes
  visible. The left column plots final values of
  \cref{eq:plasticity3d-discrete} against free DOFs; the right column plots
  the same converged energies against end-to-end wall time. All timings are
  fair 32-rank maintained runs.}
  \label{fig:plasticity3d-degree-energy}
\end{figure}

The zoom row in \Cref{fig:plasticity3d-degree-energy} is essential because the
full figure is dominated by coarse $P_1$ variation. Under the corrected
glued-bottom constraint, the first matched comparison is
$P_4(L_1)$ versus $P_1(L_3)$ at \num{610964} free DOFs. The second
matched comparison is $P_1(L_4)$, $P_2(L_3)$, and $P_4(L_2)$ at
\num{4801816} free DOFs. The scientific message is
therefore that the maintained solver reaches similar terminal energies along
both $p$- and $h$-refined paths, while higher order on coarser meshes buys
lower energy at matched DOFs only by paying a clear wall-time premium.

\begin{figure}[H]
  \centering
  \includegraphics[width=\PaperFigureMedium]{plasticity3d_recommended_scaling.pdf}
  \caption{Historical strong scaling for the promoted Plasticity3D maintained
  stack: \localpmg{} on $P_4(L_2)$ at
  $\lambda_{\mathrm{sr}} = 1.0$. This is a separate timing study from the
  glued-bottom $\lambda_{\mathrm{sr}} = 1.55$ results shown in
  \cref{fig:plasticity3d-state,fig:plasticity3d-convergence,fig:plasticity3d-degree-energy}.}
  \label{fig:plasticity3d-scaling}
\end{figure}

\begin{table}[H]
  \centering
  \begin{tabular*}{\textwidth}{@{}c@{\extracolsep{\fill}} c c c c c c@{}}
\toprule
Ranks & Wall time [s] & Solve time [s] & Speedup & Efficiency & Newton iters & Krylov iters \\
\midrule
1 & \num{5424} & \num{5176} & \num{1} & \num{1} & 17 & 114 \\
2 & \num{2914} & \num{2777} & \num{1.861} & \num{0.931} & 17 & 114 \\
4 & \num{1717} & \num{1634} & \num{3.16} & \num{0.79} & 17 & 114 \\
8 & \num{1164} & \num{1107} & \num{4.658} & \num{0.582} & 17 & 114 \\
16 & \num{608} & \num{575} & \num{8.916} & \num{0.557} & 17 & 114 \\
32 & \num{300} & \num{284} & \num{18.053} & \num{0.564} & 17 & 114 \\
\bottomrule
\end{tabular*}

  \caption{Historical Plasticity3D strong scaling for \localpmg{} on
  $P_4(L_2)$ at $\lambda_{\mathrm{sr}} = 1.0$. These rows are useful as
  solver evidence, but they should not be read as a direct timing continuation
  of the glued-bottom $\lambda_{\mathrm{sr}} = 1.55$ benchmark figures.}
  \label{tab:plasticity3d-scaling}
\end{table}

\Cref{fig:plasticity3d-scaling,tab:plasticity3d-scaling} give the paper's
largest maintained distributed mechanics timing result. This scaling evidence uses
the promoted $P_4(L_2)$ mesh family but a different load factor,
$\lambda_{\mathrm{sr}} = 1.0$, so it complements rather than extends the
glued-bottom $\lambda_{\mathrm{sr}} = 1.55$ benchmark figures above. Within
that timing study, the reported nonlinear and linear iteration counts
remain matched on 1/2/4/8/16/32 ranks, while the end-to-end time drops from
\SI{5424}{s} on one rank to \SI{300}{s} on 32 ranks. Additional
source-vs-local and continuation evidence is reported later in
\Cref{sec:solver-evidence}.

A final fixed-work diagnostic repeats the glued-bottom
$P_4(L_2)$, $\lambda_{\mathrm{sr}}=\num{1.55}$ case with the current
constitutive-AD local assembly and MUMPS-backed PMG stack. The PMG hierarchy is
the same-mesh $P_4 \rightarrow P_2 \rightarrow P_1$ hierarchy, the nonlinear
method is Armijo Newton, and the Krylov tolerance is \num{1e-1}. Each run is
intentionally capped after five Newton iterations, so the raw solver status is
``maximum number of iterations reached'' even though the fixed work is matched.
\Cref{fig:plasticity3d-local-karolina-scaling,tab:plasticity3d-local-karolina-scaling}
overlay the local workstation sweep with the Karolina CPU-node sweep at
16 MPI ranks per node. All plotted rows use 29 Krylov iterations and reach the
same capped-state energy and gradient norm to the precision shown in the table;
the figure is therefore a fixed-work implementation diagnostic rather than a
new convergence claim.

\begin{figure}[H]
  \centering
  \includegraphics[width=\PaperFigureMedium]{plasticity3d_local_vs_karolina_scaling.pdf}
  \caption{Local and Karolina fixed-work scaling for the glued-bottom
  Plasticity3D $P_4(L_2)$ case at $\lambda_{\mathrm{sr}}=\num{1.55}$. Dashed
  references are ideal $1/r$ lines anchored separately to the first point of
  each platform.}
  \label{fig:plasticity3d-local-karolina-scaling}
\end{figure}

\begin{table}[H]
  \centering
  \begin{tabular*}{\textwidth}{@{}l@{\extracolsep{\fill}} c c c c c c c@{}}
\toprule
Platform & Ranks & Nodes & Solver total [s] & Solve [s] & Krylov iters & Energy & $\|g\|$ \\
\midrule
local workstation & 1 & -- & \num{1626} & \num{1370} & 29 & \num{-3013147} & \num{488.443} \\
local workstation & 2 & -- & \num{884} & \num{733} & 29 & \num{-3013147} & \num{488.443} \\
local workstation & 4 & -- & \num{518} & \num{420} & 29 & \num{-3013147} & \num{488.443} \\
local workstation & 8 & -- & \num{347} & \num{298} & 29 & \num{-3013147} & \num{488.443} \\
local workstation & 16 & -- & \num{181} & \num{155} & 29 & \num{-3013147} & \num{488.443} \\
local workstation & 32 & -- & \num{106} & \num{90.9} & 29 & \num{-3013147} & \num{488.443} \\
Karolina & 16 & 1 & \num{187} & \num{147} & 29 & \num{-3013147} & \num{488.443} \\
Karolina & 32 & 2 & \num{112} & \num{88.8} & 29 & \num{-3013147} & \num{488.443} \\
Karolina & 64 & 4 & \num{74.5} & \num{60} & 29 & \num{-3013147} & \num{488.443} \\
Karolina & 128 & 8 & \num{55.2} & \num{45.2} & 29 & \num{-3013147} & \num{488.443} \\
\bottomrule
\end{tabular*}

  \caption{Timings for the capped five-Newton-iteration Plasticity3D
  diagnostic in \cref{fig:plasticity3d-local-karolina-scaling}. Local rows are
  single-workstation MPI runs; Karolina rows use 16 ranks per CPU node.}
  \label{tab:plasticity3d-local-karolina-scaling}
\end{table}

\subsection{Topology Optimization}
\label{sec:num-topology}

The reduced design problem, mechanics subproblem, and objective history are
defined in
\Cref{sec:bench-topology,eq:topology-mechanics,eq:topology-reduced-objective}.
The numerical question here is distributed outer-loop stability: the mechanics
solve and the design update must remain synchronized under continuation in $p$
and move-regularized updates. Because the outer-loop work can vary with the
continuation schedule, the timing evidence here should be read as end-to-end
workflow scaling rather than fixed-work strong scaling.

\begin{figure}[H]
  \centering
  \includegraphics[width=\PaperFigureMedium]{topology_scaling.pdf}
  \caption{Topology end-to-end parallel timing trend for the maintained
  fine-grid run. The plotted wall times correspond to the coupled reduced-design
  workflow defined by \cref{eq:topology-reduced-objective}.}
  \label{fig:topology-results}
\end{figure}

\begin{table}[H]
  \centering
  \begin{tabular*}{\textwidth}{@{}c@{\extracolsep{\fill}} c c c c c c@{}}
\toprule
Ranks & Wall time [s] & Solve time [s] & Outer iters & $p$ & Compliance & Volume fraction \\
\midrule
1 & \num{208} & \num{200} & 65 & \num{4.2} & \num{9.1559} & \num{0.3884} \\
2 & \num{138} & \num{134} & 72 & \num{5.6} & \num{8.9473} & \num{0.3932} \\
4 & \num{95.7} & \num{92.8} & 69 & \num{4.6} & \num{9.1684} & \num{0.385} \\
8 & \num{82.1} & \num{80.1} & 67 & \num{7} & \num{9.2178} & \num{0.3932} \\
16 & \num{43.3} & \num{41.9} & 66 & \num{5.8} & \num{9.6859} & \num{0.3798} \\
32 & \num{25.1} & \num{23.7} & 66 & \num{6.6} & \num{9.9074} & \num{0.3749} \\
\bottomrule
\end{tabular*}

  \caption{Parallel topology summary for the maintained fine-grid study,
  reported for the reduced objective in \cref{eq:topology-reduced-objective}.}
  \label{tab:topology-summary}
\end{table}

The topology benchmark completes the paper's computational picture. Here the
mainline \jaxpetsc{} path is less about exact Hessians than about stable
distributed design-mechanics coupling, while \purejax{} remains the compact
serial reference path discussed earlier in \Cref{tab:reference-availability}.
